\chapter{基于截断分位数算法的机械臂操作任务专家策略训练}\label{chap:reinforce}
\section{引言}\label{sec:reinforce-intro}
在\cref{chap:lead}中提到，基于LLMs的机械臂控制方法在实现底层动作对齐以及指令微调时都需要高质量的专家示范数据作为支撑。传统获取专家数据的方式往往需要借助人工操作示教或者硬编码方法，人工示教的成本偏高而且受到示教者自身经验水平的制约，硬编码方法则在灵活度上难以满足动态接触类操作任务的需求。基于上述原因，本章运用强化学习方法为各项基础操作任务分别训练了单任务专家模型，以此来为后续数据集构建工作提供数据来源。

本章以MuJoCo物理仿真引擎为基础构建Gymnasium-Robotics仿真环境，围绕Fetch机械臂的Reach、Push、Slide以及PickAndPlace四项基础操作任务展开研究，具体组织如下。\cref{sec:rl-theory}概述强化学习基础理论以及本章所选用的截断分位数(TQC)算法与后验经验回放(HER)机制。\cref{sec:rl-environment}针对四项任务各自的物理特性，分别设计了对应的状态空间、动作空间以及奖励函数。\cref{sec:expert-training}凭借TQC以及HER算法对四项任务开展独立的策略训练，并且借助消融实验来验证HER所起到的作用，最终对各专家模型的评估成功率进行定量分析。

\section{强化学习基础理论和算法}\label{sec:rl-theory}
\subsection{强化学习基础理论}\label{sec:rl-basics}
强化学习通过智能体同环境的持续交互，以试错方式最大化累积奖励，以进行策略优化。在机械臂控制等序贯决策任务中，由于未来状态仅依赖于当前状态与执行动作，该过程满足马尔可夫性质，可以被建模为马尔可夫决策过程(MDP)。

MDP定义为一个五元组$\langle \mathcal{S}, \mathcal{A}, P, r, \gamma \rangle$，其中$\mathcal{S}$代表智能体自身与环境所有可观测状态的集合；$s$表示当前时刻的状态；$\mathcal{A}$是智能体所能执行的动作空间；$a$表示当前时刻下智能体执行的动作；$P$是下一时间步状态$s'$与当前时刻状态$s$及执行动作$a$间的状态转移概率分布；$r$代表奖励函数，是基于当前状态$s$执行动作$a$后，在下一时刻下环境给出的奖励；$\gamma \in [0,1)$是折扣因子，反映未来奖励的重要性。由此，以当前时刻$t$开始计算的累计折扣回报$G_t$可以表示为：
\begin{align}
    G_t &= \sum_{k=0}^{\infty} \gamma^k r_{t+k+1}
\end{align}

强化学习的目标是学习策略$\pi(a|s)$以最大化期望累计折扣回报，其中$\pi(a|s)$表示在状态$s$下选择动作$a$的概率。为评估策略优劣，定义策略$\pi$下的状态价值函数与动作价值函数如下：
\begin{align}
    V^\pi(s) &= \mathbb{E}_\pi[G_t | s_t = s] \\
    Q^\pi(s, a) &= \mathbb{E}_\pi[G_t | s_t = s, a_t = a]
\end{align}
基于上式给出的定义，记当前时间步$t$下的状态、动作和环境奖励为$(s, a, r)$，下一时刻的对应值为$(s', a', r')$。状态与动作价值函数递推关系可以用贝尔曼方程表示：
\begin{align}
    V^\pi(s) &= r' + \gamma \sum_{a} \pi(a|s) \sum_{s'} P(s'|s,a) V^\pi(s') \\
    Q^\pi(s, a) &=r' + \gamma \sum_{s'} P(s'|s,a) \sum_{a'} \pi(a'|s') Q^\pi(s', a')
\end{align}
强化学习的目标是寻找最优策略$\pi^*$，使得在所有状态下的价值函数达到最大：
\begin{align}
    V^{\pi^*}(s) &= r' + \gamma \max_{a \in \mathcal{A}}{\sum_{s'} P(s'|s,a) V^{\pi^*}(s')} \\
    Q^{\pi^*}(s, a) &=r' + \gamma \max_{a \in \mathcal{A}}{\sum_{s'} P(s'|s,a) \sum_{a'} Q^{\pi^*}(s', a')}
\end{align}

主流深度强化学习算法凭借神经网络近似上述策略函数$\pi(a|s)$和价值函数$V(s)$、$Q(s,a)$，让它能够应用于高维连续状态以及动作空间的机械臂任务。

\subsection{截断分位数算法}\label{sec:tqc}
当前强化学习算法按最优化目标主要分为三类，基于值函数的方法、基于策略的方法以及Actor-Critic的方法。基于值函数的方法侧重于近似状态价值函数$V(s)$或动作价值函数$Q(s,a)$，并基于价值函数使用贪心算法得到行为策略，典型算法包括Q-Learning\cite{watkins1992q}、DQN\cite{mnih2013playing}等。这类方法维护一个离散值函数表，对连续动作空间的处理能力有限。基于策略的方法直接参数化策略$\pi_\theta(a|s)$，凭借梯度下降等方式最大化期望回报来优化策略参数$\theta$，如REINFORCE\cite{williams1992simple}、PPO\cite{schulman2017proximal}等。这类方法适用于连续动作空间，不过存在采样轨迹数据方差过大所导致的训练不稳定等一系列问题。演员-评论家方法同时维护Actor策略网络与Critic价值网络，既能处理连续动作空间，又具有较好的训练稳定性，如SAC\cite{haarnoja2018soft}、DDPG\cite{lillicrap2020continuous}等。本章则使用基于演员-评论家方法的截断分位数(TQC)算法\cite{kuznetsov2020tqc}来训练机械臂操作任务的专家模型。

为缓解传统Actor-Critic架构中单一标量评估导致的过度估计偏差，TQC采用分布式强化学习的方法，借助多网络集成的分位数回归来近似回报的概率分布。TQC首先把累积回报$G_t$明确表示为一个同状态$s$以及动作$a$相关的随机变量$G(s,a)$，其期望即为传统的动作价值函数，即$Q^\pi(s,a) = \mathbb{E}_\pi[G(s,a)]$，定义该累积回报分布为$\mathcal{Z}(s,a)$。为了明确其概率分布，TQC定义了$N$个Critic网络$\{\theta_{\psi_i}(s,a)\}_{i=1}^N$，每个Critic网络输出的对应离散累积回报分布$z_i(s,a)$则是由$M$个分位数值$\{\theta_{\psi_i}^{j}(s,a)\}_{j=1}^M$进行近似获得。因此，对于任意给定的状态动作对$(s,a)$，总Critic网络能够输出$N \times M$个分位数值。

由于实际上智能体并不能获取环境的状态$s$，只能基于自己的观测$o$来执行动作，其中存在转移函数$o = f(s)$。作为一种离线策略算法，TQC把各个Critic网络的在线参数$\psi_n$同目标参数$\bar{\psi_n}$解耦，因此需要构建经验回放池$\mathcal{D}$来存储基于在线网络运行的轨迹，该轨迹被记录为四元组$(o, a, r, o')$进行存储。在线网络凭借经验回放池中的轨迹数据进行梯度更新，而目标网络则凭借软更新系数$\tau$对在线网络参数进行指数平滑跟踪，以提高训练数据利用率。

为了将离散的分度值拟合为累积回报分布，TQC引入了Dirac测度$\delta_x$，定义为：
\begin{align}
    \delta_x(A)=\left\{\begin{matrix}
        1, & x \in A\\
        0, & x \notin A
        \end{matrix}\right.
\end{align}
由此给每个在合理区间范围$A$内的分位度值$\theta_{\psi_i}^{j}(o,a)$分配同等概率，构造Critic网络$\theta_{\psi_i}$的离散累积回报分布：
\begin{align}
    z_{\psi_i}(o,a) := \frac{1}{M} \sum_{j=1}^M \delta_{\theta_{\psi_i}^{j}(o,a)}(A)
\end{align}
再由此搜集所有Critic网络输出的$N \times M$预测分度值，构建集合进行池化：
\begin{align}
    \mathcal{Z}(o,a) := \{z_{\psi_i}(o,a)|i=1,\ldots,N\} &= \{\theta_{\psi_i}^{j}(o,a)|i=1,\ldots,N; j=1,\ldots,M\}
\end{align}
为控制过度估计偏差，TQC算法在目标分布的计算中引入截断机制，去除数值过大的估计。把池化集合$\mathcal{Z}(o,a)$按升序排列，记排序后的第$k$个元素为$\theta_{k}(o,a)$。基于池化后的分布根据预设截断数量$d$从排序后的集合中强制舍弃数值最大的$d \times N$个分位数，仅保留较小的$K = (M - d) \times N$个分位数，如\cref{fig:tqc-quantile}所示。经过截断后的最新分布如下：
\begin{align}
    \mathcal{Z}(o,a) := \frac{1}{K} \sum_{k=1}^K \delta_{\theta_{k}(o,a)}(A)
\end{align}

之后，引入SAC\cite{haarnoja2018soft}所使用的最大熵框架来鼓励探索：高熵时偏向探索新动作，低熵时偏向运用当前最优策略，从而降低收敛到局部最优解的概率。在这一框架下，结合目标策略$\pi$的熵定义为：
\begin{align}
    \mathcal{H}(\pi(\cdot|o)) = -\sum_{a} \pi(a|o) \log \pi(a|o)
\end{align}
对应要优化的目标函数也定义为最大化累积回报与熵的和，给定对应温度系数$\alpha$，优化目标可表示为：
\begin{align}
    J(\pi) = \max_a \mathbb{E}_{\pi(\cdot|o)}\left[ \mathcal{Z}^\pi(o,a) + \alpha \mathcal{H}(\pi(\cdot|o)) \right]
\end{align}
对应的熵正则化项为$\alpha \log \pi(a|o)$，构建对应目标的原子分布$\{y_k(o|a)\}_{k=1}^K$：
\begin{align}
    y_k(o, a) = r(o, a) + \gamma \left[ z_{k}(o', a') - \alpha \log \pi(a'|o') \right]
\end{align}
因此最终构建的目标分布可以表示为：
\begin{align}
    \mathcal{Y}(o, a) := \frac{1}{K} \sum_{k=1}^K \delta_{y_k(o,a)}(A)
\end{align}
在实际训练中，基于智能体同环境的交互进行采样得到目标分布的估计值$\hat{\mathcal{Y}}$，并且运用分位数回归损失函数来更新Critic网络参数$\theta_{\psi_n}$，同时运用策略梯度方法来更新Actor网络参数$\phi$，以最大化目标函数$J(\pi)$。给定分位数$\tau$的Huber回归损失函数为：
\begin{align}
\rho_\tau^H(u) = |\tau - \mathbb{I}_{{u < 0}}| \cdot \mathcal{L}_H^1(u)
\end{align}
其中Huber损失函数$\mathcal{L}_H^1(u)$在误差较小时表现为二次损失，在误差较大时表现为线性损失，从而兼顾了均方误差的稳定性和绝对误差的鲁棒性。在对应公式中定义为：
\begin{align}
    \mathcal{L}_H^1(u) = \left\{\begin{matrix}
        \frac{1}{2}u^2, & |u| \leq 1\\
        |u| - \frac{1}{2}, & |u| > 1
        \end{matrix}\right.
\end{align}
记$k=\frac{K}{N}$，由此可以给出Critic网络的损失函数及优化目标：
\begin{align}
    \mathcal{L}_k(o,a;\psi_n) &= \frac{1}{KM} \sum_{m=1}^M \sum_{i=1}^K \rho_{\tau_m}^H(y_i(o,a) - \theta_{\psi_n}^{m}(o,a)) \\ 
    J_Z(\psi_n)&=\mathbb{E}_{\mathcal{D},\pi}[\mathcal{L}_k(o,a;\psi_n)]\label{eq:critic_loss}
\end{align}
再基于分布提取总Critic网络对动作状态价值函数的估计值：
\begin{align}
    \hat{Q}_{TQC}^{\pi}(o,a) = \frac{1}{K} \sum_{i=1}^K \theta_{i}(o,a)
\end{align}
来给出策略Actor网络的优化目标：
\begin{align}\label{eq:actor_loss}
    J_{\pi}(\phi) = -\mathbb{E}_{\mathcal{D}} \left[ \mathbb{E}_{\pi_\phi(\cdot|o)} \left[ \frac{1}{K} \sum_{i=1}^K \theta_{i}(o,a) - \alpha \log \pi_\phi(a|o) \right] \right]
\end{align}
其中，熵的温度系数$\alpha$也可以通过引入一个单独的优化目标来进行自动调整，以确保策略在训练过程中保持适当的探索程度。对应的温度系数优化目标可以定义为：
\begin{align}\label{eq:alpha_loss}
    J(\alpha) = \mathbb{E}_{a \sim \pi_\phi(\cdot|o)}\left[ -\alpha \log \pi_\phi(a|o) - \alpha \mathcal{H}_T \right]
\end{align}

据此，TQC算法的完整流程如\cref{algo:tqc_flow}所示。

\begin{algorithm}
  \caption{TQC 算法流程}\label{algo:tqc_flow}
  \begin{algorithmic}[1]
    \Require 随机初始化的策略网络参数$\phi$，$N$个Critic网络参数$\psi_n$及其目标网络$\bar{\psi}_n$, $n \in [1..N]$
    \Require 经验回放池$\mathcal{D} = \emptyset$，目标熵$\mathcal{H}_T$，温度系数 $\alpha$，软更新系数$\tau$

        \For{回合$e=1 \to E$}
        \State 重置环境状态$s_0$
        \For{每一个环境采集步$t=0 \to T-1$}
            \State 使用当前策略 $\pi_\phi$ 采集样本 $(s_t, a_t, r_t, s_{t+1})$
            \State 将样本存入回放池：$\mathcal{D} \leftarrow \mathcal{D} \cup \{(s_t, a_t, r_t, s_{t+1})\}$
        \EndFor
        
        \For{每一个梯度更新步}
            \State 从回放池$\mathcal{D}$中随机采样一个批次的数据
            \State 更新温度系数$\alpha$：
            \State \quad $\alpha \leftarrow \alpha - \lambda_\alpha \hat{\nabla}_\alpha J(\alpha)$ \Comment{源自式\eqref{eq:alpha_loss}}
            
            \State 更新策略网络Actor参数$\phi$：
            \State \quad $\phi \leftarrow \phi - \lambda_\pi \hat{\nabla}_\phi J_\pi(\phi)$ \Comment{源自式\eqref{eq:actor_loss}}
            
            \State 更新各在线分位数Critic网络参数$\psi_n$：
            \For{$n = 1$ \textbf{to} $N$}
                \State $\psi_n \leftarrow \psi_n - \lambda_Z \hat{\nabla}_{\psi_n} J_Z(\psi_n)$ \Comment{源自式\eqref{eq:critic_loss}}
            \EndFor
            
            \State 软更新目标Critic网络参数$\bar{\psi}_n$：
            \State \quad $\bar{\psi}_n \leftarrow \tau \psi_n + (1 - \tau) \bar{\psi}_n, \quad n \in [1..N]$
        \EndFor
    \EndFor
    \State \Return 最终策略$\pi_\phi$与Critics$Z_{\psi_n}$
  \end{algorithmic}
\end{algorithm}

\subsection{后验经验回放机制}\label{sec:her}
在上节\cref{algo:tqc_flow}中，进行梯度更新的数据来源于智能体在探索过程中记录轨迹的经验回放池$\mathcal{D}$。在本章所使用的机械臂操作任务环境中，由于稀疏奖励特性，智能体在初始阶段很难获得正向奖励信号，训练效率低。为此，本节引入了后验经验回放(HER)\cite{andrychowicz2018hindsight}机制。HER把智能体在失败轨迹中实际达到的状态作为新的目标状态进行重新标记，来重构经验回放池$\mathcal{D}$，使原本失败的轨迹也能产生学习信号，加速策略收敛。

\cref{sec:tqc}中定义了经验回放池$\mathcal{D}$中记录的四元组$(o, a, r, o')$。由于环境实际给予的奖励同目标状态$g$相关，即$r=R(s, g)$，故可以额外增加一维度$g$记录进入轨迹中，构成五元组$(o, a, r, o', g)$。HER通过在采样阶段对目标g重新标记来重构转移元组。给定一个由环境收集的完整回合轨迹序列$\tau = \{(o_0, a_0), (o_1, a_1), \ldots, (o_T, a_T)\}$，对于时间步$t$的转移记录，HER对目标g进行了重构。从同一回合中时间步$t' > t$的实际物理状态中提取特征，映射为一个新的虚拟目标状态$g'$。给定观测$o$和目标状态$g$间的映射函数$\phi$，可以将新目标状态表示为$g' = \phi(o_{t'})$，并基于该新目标状态重新计算奖励$r' = R(s, g')$。因此，原始的转移记录$(o, a, r, o', g)$被重构为新的转移记录$(o, a, r', o', g')$，并存储回经验回放池$\mathcal{D}$中。

\section{机械臂操作任务建模}\label{sec:rl-environment}
本节运用基于MuJoCo物理引擎的Gymnasium-Robotics作为基准测试环境，选取以7自由度Fetch移动机械臂为核心的Reach、Push、Slide以及PickAndPlace四项基础操作任务。在\cref{sec:her}所介绍的目标条件马尔可夫决策过程框架下，本节将把这四项任务的物理状态、控制指令以及反馈机制统一建模为状态与目标空间、动作空间以及奖励空间。环境建模的详细信息见\cref{chap:detail-env}。

\subsection{状态与目标空间建模}\label{sec:state-space}
在状态与目标空间方面，Gymnasium-Robotics环境下Fetch任务的观测空间的返回值是一个字典，其中包含三个键值：当前观测$o_t$、当前达成目标$g_a$以及期望目标$g$。依据是否涉及被操作物体，这四项任务的观测维度存在着差异。

Reach任务要求把机械臂末端执行器移动至空间指定位置，不涉及与外部物体的接触和交互，因此它的观测$o_t$仅包含机械臂自身的运动学信息，也就是末端执行器相对于全局坐标系的坐标与速度以及夹爪左右指节的相对位置和速度，构成一个10维连续向量。其目标空间$G$为3维向量，表示末端执行器需要到达的目标空间坐标$(x, y, z)$。

Push、Slide以及PickAndPlace任务则要求机械臂同外部物体进行交互。这几项任务的观测$o_t$扩展为25维连续向量，除了包含上述10维的机械臂本体状态之外，还引入了物体的三维空间位置、欧拉角、线速度、角速度以及物体相对于末端执行器的相对位移与相对速度等参数作为新的观测信息。其目标空间$G$同样为3维向量，不过所指的是目标物体需要到达的期望空间坐标。

\subsection{动作空间建模}\label{sec:action-space}
四项任务的底层控制指令统一建模为4维的连续动作空间$\mathcal{A}$，其取值范围为$\mathcal{A} = [-1.0, 1.0]^4$。该动作向量在四个任务下的定义相同，不过存在部分维度在特定环境下不被激活的情况。

动作空间前三个维度$(a_0, a_1, a_2)$代表机械臂末端执行器在全局笛卡尔坐标系下$(x, y, z)$轴方向上的期望相对位移量。通过MuJoCo的逆运动学求解器，结合机械臂自身的固定特性，这3维笛卡尔位移可被映射为Fetch机械臂7个关节的控制扭矩。

第四个维度$a$用于控制夹爪的开合状态。在Reach、Push以及Slide任务中，由于无需抓取，该维度参量不会激活对应的动力学运算模块来实现夹爪控制。而在PickAndPlace任务中，该维度则会控制两侧指尖的相对距离，实际运行中正数代表张开，负数代表闭合。

\subsection{奖励空间建模}\label{sec:reward-space}
针对这四项机械臂操作任务，本研究统一运用稀疏奖励机制，以此来避免主观设计稠密奖励所引入的局部最优陷阱。定义距离函数$d(g_a, g)$为当前达成目标同期望目标在三维空间中的距离，给定距离容差阈值$\epsilon=0.05$米，对应时间步$t$的即时奖励$r_t$定义为阶跃函数：
\begin{align}
    r_t = 
    \begin{cases}
        0, & \text{if } d(g_a, g) \leq \epsilon \\
        -1, & \text{otherwise}
    \end{cases}
\end{align}

也就是说，智能体在绝大多数交互步骤中都会收到-1惩罚，仅在达成目标时获得奖励。这种稀疏奖励机制加大了初始阶段的探索难度，不过在结合TQC与HER方法之后能够有效训练出目标导向的策略，避免策略网络陷于局部最优。

\section{专家模型训练与性能评估}\label{sec:expert-training}
\subsection{实验设置与网络架构}\label{sec:experiment-setup}
由于各项任务的复杂度不同，本节针对每项任务的特性分别设计了不同的网络架构以及超参数配置。Gymnasium-Robotics环境返回的观测数据为包含观测与目标的字典。本研究在策略网络和价值网络上运用了多输入的架构来处理该字典。该架构先在网络输入层借助特征提取器把当前状态$o_t$与目标状态$g$进行拼接融合，再输入至由简单全连接神经网络构成的多层感知机(MLP)当中。

考虑到各项任务在状态空间维度以及接触动力学方面存在着区别，本研究对网络复杂度以及训练周期进行了针对性设计。并结合了\cref{sec:tqc}中介绍的TQC算法以及\cref{sec:her}介绍的HER方法来处理稀疏奖励问题。

涉及物体交互的任务在观测信息复杂度以及状态空间维度上要求更高。Push、PickAndPlace以及Slide任务涉及25维的复杂状态空间，并且要求机械臂同物体发生持续或瞬间碰撞。这三项任务选用了$[512, 512, 512]$维度MLP架构作为演员和评论家网络的基础架构，实践中设定TQC算法中评论家网络集成数量$N=2$来输出合适的分位数分布，每个网络输出的分位数数量为默认值25。

Push以及PickAndPlace任务的训练时间步设定为$1 \times 10^6$，由于面对稀疏奖励环境，折扣因子设定为$\gamma = 0.98$，使模型更重视远期奖励。

Slide任务包含远距离和近距离的目标，要求机械臂既能通过稳健的推移操作将目标物体移动到近距离目标，又能凭借精确的击打使物体滑动到远距离目标。出于任务复杂度方面的考虑，该任务的训练步数延长至$3 \times 10^6$步。

Reach任务仅涉及10维的自身运动学状态，控制逻辑相对简单。为避免过度参数化导致的过拟合以及计算资源浪费，该任务运用$[64, 64]$维的轻量化MLP网络架构，并且仅保留$N=1$个Critic网络，网络输出的分位数数量为默认值25。在此基础上，为了加速初期训练的梯度平稳性，对Reach环境开启了状态观测值的归一化，借助动态维护状态的历史均值与方差把输入神经网络的特征缩放至标准正态分布。该任务较为简单而且完成所需步数少，训练时间步压缩为$2 \times 10^4$步，折扣因子下调为$\gamma = 0.95$，以此来聚焦于更短期的运动轨迹。

四项任务均设定HER重标记比例参数$\beta = 4$，也就是说对于每1条从环境中采集的真实转移轨迹，HER机制将在同一回合的后续时间步中提取状态作为虚拟目标，额外生成4条虚拟的成功经验。

实际运用的网络以及训练参数如下\cref{tab:fetch-hyperparams}所示。

\begin{table}[!htb]
  \caption{Fetch类任务训练超参数配置表}\label{tab:fetch-hyperparams}
  \centering
  \small
  \begin{tabular}{@{} l ccc @{}} \toprule
    \textbf{参数名称} & \textbf{Push / PickAndPlace} & \textbf{Slide} & \textbf{Reach} \\ \midrule
    隐藏层结构 & [512, 512, 512] & [512, 512, 512] & [64, 64] \\
    Critic网络数量$N$ & 2 & 2 & 1 \\
    总训练步数 & $1 \times 10^6$ & $3 \times 10^6$ & $2 \times 10^4$ \\
    折扣因子$\gamma$ & 0.98 & 0.98 & 0.95 \\
    学习率 & $1 \times 10^{-3}$ & $1 \times 10^{-3}$ & $1 \times 10^{-3}$ \\
    经验池$\mathcal{D}$容量 & $1 \times 10^6$ & $1 \times 10^6$ & $1 \times 10^6$ \\
    每次更新采样批次大小 & 512 & 512 & 256 \\
    软更新系数$\tau$ & 0.005 & 0.005 & 0.005 \\
    虚拟目标采样数$\beta$ & 4 & 4 & 4 \\
    初始状态观测归一化 & False & False & True \\ \bottomrule
  \end{tabular}
\end{table}

\subsection{训练过程与收敛性分析}\label{sec:convergence-analysis}
为评估单任务专家模型的学习效率以及策略稳定性，本节使用了周期性在线评估机制。即训练过程中，模型每隔固定时间步在独立测试环境中执行50个评估回合，策略网络不进行探索而直接输出当前最优动作，并把评估成功率记录至TensorBoard日志。为了验证HER机制对训练效率的影响，本研究同时记录了相同超参数下不使用HER的TQC基线训练过程。训练曲线运用窗口长度为数据总量10\%的三阶Savitzky-Golay滤波器进行平滑处理，阴影区域由滤波器缓冲区内10至90百分位数的统计范围构成，用于表征训练过程中的波动程度。四项任务在TQC+HER与纯TQC两种配置下的训练收敛曲线如\cref{fig:tqc_trainline}所示。

\begin{figure}[!htb]
  \centering
  \begin{subfigure}[b]{0.48\textwidth}
    \centering
    \includegraphics[width=\textwidth]{figures/FetchReach.png}
    \caption{FetchReach}\label{fig:tqc-reach}
  \end{subfigure}
  \hfill
  \begin{subfigure}[b]{0.48\textwidth}
    \centering
    \includegraphics[width=\textwidth]{figures/FetchPush.png}
    \caption{FetchPush}\label{fig:tqc-push}
  \end{subfigure}
  \vspace{2pt}
  \begin{subfigure}[b]{0.48\textwidth}
    \centering
    \includegraphics[width=\textwidth]{figures/FetchSlide.png}
    \caption{FetchSlide}\label{fig:tqc-slide}
  \end{subfigure}
  \hfill
  \begin{subfigure}[b]{0.48\textwidth}
    \centering
    \includegraphics[width=\textwidth]{figures/FetchPickAndPlace.png}
    \caption{FetchPickAndPlace}\label{fig:tqc-pick}
  \end{subfigure}
  \caption{四项Fetch任务在TQC与TQC+HER配置下的训练收敛曲线对比。实线为Savitzky-Golay滤波后的中心趋势，阴影区域表示第10至第90百分位的统计波动范围。}\label{fig:tqc_trainline}
\end{figure}

\cref{fig:tqc_trainline}所示的四组收敛曲线展示了HER机制在不同任务复杂度下对训练过程的影响。Reach任务的10维状态空间结构简单，TQC与TQC+HER两种方法均成功收敛至接近100\%的成功率，其中TQC+HER在约6400步突破80\%成功率，纯TQC则需要约9800步，提升了约1.5倍。在简单任务中，HER主要凭借加速初期探索来缩短训练周期，并非收敛的必要条件。

Push任务中，TQC+HER在221K步突破80\%成功率并且最终稳定在97.9\%，纯TQC虽最终也达到97.5\%的成功率，不过达到80\%成功率所需的步数约为378K步，相比TQC+HER慢了约1.7倍。PickAndPlace任务中两组配置的则在最终收敛策略的成功率上存在明显差异，TQC+HER在约492K步达到80\%以上成功率并且最终稳定在99.0\%，而纯TQC在整个$1 \times 10^6$步训练周期内始终未突破80\%，最终成功率仅为56.7\%。从最终收敛状态来看，TQC+HER的策略波动范围小于纯TQC，说明引入HER能够提升训练过程中的策略稳定性。

Slide任务呈现出了同其他三项任务不同的趋势。TQC+HER在约279K步即突破50\%成功率，初期上升速度快于纯TQC的664K步，不过在训练中后期逐渐趋于饱和，最终成功率为80.8\%。相比之下，纯TQC虽然初期学习较慢，但最终成功率达到了90.6\%，明显高于TQC+HER。从动力学角度分析，Slide任务要求机械臂借助精确的碰撞冲量把物块击打至远距离目标位置，目标分布涵盖近距与远距目标。HER的future目标重标记策略把实际到达的物体位置作为虚拟目标，在Slide任务中可能引入过多同原始目标分布不一致的偏向性监督信号，使策略倾向于学习将物块推至任意中间位置而非针对远距离目标的精确击打策略，导致策略在真实评估目标上的泛化性能下降。基于此消融实验结果，本研究最终选取纯TQC模型作为Slide任务的专家模型。

训练完成后，各任务专家模型在独立评估中的成功率如\cref{tab:expert-success-rate}所示。其中Reach、Push以及PickAndPlace任务运用TQC+HER配置，Slide任务运用纯TQC配置。Reach以及PickAndPlace任务的成功率分别为100.0\%和99.0\%。Push任务成功率为97.9\%。Slide任务成功率为90.6\%，在四项任务中最低。四项任务对应专家策略的成功率均超过90\%，为后续\cref{chap:dataset}当中多任务微调数据集的构建提供了轨迹数据来源。

\begin{table}[htbp]
\centering
\caption{各任务专家模型评估成功率}\label{tab:expert-success-rate}
\begin{tabular}{lcccc}
\toprule
任务 & Reach & Push & PickAndPlace & Slide \\
\midrule
所选配置 & TQC+HER & TQC+HER & TQC+HER & TQC \\
评估成功率 (\%) & 100.0 & 97.9 & 99.0 & 90.6 \\
\bottomrule
\end{tabular}
\end{table}

\subsection{小结}
本章基于TQC算法以及HER机制，针对四项Fetch机械臂操作任务分别训练了单任务专家策略。消融实验结果表明，HER机制在稀疏奖励环境中普遍起到了加速训练收敛的作用，其效果随任务复杂度的提升而增强。不过，对于Slide等目标空间分布同实际状态分布存在较大差异的任务，HER的虚拟目标重标记可能引入过多的偏向性监督信号，反而降低最终策略质量。本研究借助对比消融实验为每项任务选择了最优配置，训练后的专家模型在四项任务上达成了\cref{tab:expert-success-rate}所示的成功率，为\cref{chap:dataset}当中多任务微调数据集的构建提供了高质量的轨迹数据来源。
